\begin{tikzpicture}
%HB9ARK 16-NOV-2026
\documentclass[tikz,border=5pt]{standalone}
\usetikzlibrary{decorations.pathreplacing}

\begin{document}
\begin{tikzpicture}[x=1cm,y=1cm,line width=0.8pt]

\def\H{1.2}
\def\L{0}

% Pegelbezeichnungen
\node[left] at (0,1.4) {Mark};
\node[left] at (0,0.16) {Space};

% Ruhezustand
\draw (-0.8,\H) -- (0,\H);
\node[above] at (-1.11,1.75) {Ruhezustand};

% Startbit
\draw (0,\H) -- (0,\L) -- (1,\L);

% Baudot Y = 10101
\draw (1,\L) -- (1,\H) -- (2,\H); % b0=1
\draw (2,\H) -- (2,\L) -- (3,\L); % b1=0
\draw (3,\L) -- (3,\H) -- (4,\H); % b2=1
\draw (4,\H) -- (4,\L) -- (5,\L); % b3=0
\draw (5,\L) -- (5,\H) -- (6,\H); % b4=1

% Stopbits (1.5 Bitzeiten)
\draw (6,\H) -- (7.5,\H);

% Bitbezeichnungen
\foreach \x/\b in {1.5/0,2.5/1,3.5/2,4.5/3,5.5/4}
  \node[below] at (\x,-0.15) {$d_{\b}$};

% Zeichenbezeichnung
\node[above] at (3.75,1.75) {Zeichen "Y" im Baudot Code  1  0  1  0  1 };

% Datenklammer
\draw[
  decorate,
  decoration={brace,mirror,amplitude=5pt}
]
(1,-0.55) -- (6,-0.55)
node[midway,below=8pt]
{Datenbits für Baudot "Y"};

% Start-/Stoppfeile
\draw[->] (0.5,-1.05) -- (0.5,-0.05);
\node[below] at (0.74,-1.11) {Startbit};

\draw[->] (6.75,-1.05) -- (6.75,1.15);
\node[below] at (6.75,-1.05) {Stopbit};

\end{tikzpicture}
\end{document}


\end{tikzpicture}